\section{Overview of the Project Implementation}

In the previous section, we established the theoretical necessity of an inward-flowing Directed Acyclic Graph (DAG) to prevent software entropy. We defined the Core Domain as the isolated center of the architecture, surrounded by the chaotic, side-effect-heavy Infrastructure layer.

However, Python does not natively enforce Clean Architecture. By default, any Python file can import any other Python file, making it dangerously easy for a developer to accidentally create a cyclic dependency or bypass an interface. To enforce our topological boundaries physically, we map the abstract concentric rings of our architecture directly to the directory tree of the \texttt{semantic\_engine} repository.

\subsection*{The Physical Directory Tree}

The codebase is strictly segregated into top-level directories that represent the distinct boundaries of the DAG. Below is a simplified representation of the repository's physical architecture:

\begin{verbatim}
src/
|-- core_interfaces/       # Pure Python protocols (The Inner Ring)
|   |-- api.py
|   |-- domain.py
|   |-- ml.py
|   |-- repository.py
|   `-- uow.py
|-- infrastructure/        # I/O, Database, APIs (The Outer Ring)
|   |-- api/
|   |-- database/
|   |-- ml/
|   `-- worker/
`-- tests/                 # Isolated verification sandbox
\end{verbatim}

\subsection*{The Rules of Import}

To maintain the acyclic nature of the system, we establish strict, non-negotiable compile-time rules for how these folders interact via Python \texttt{import} statements:

\begin{enumerate}
    \item \textbf{The Core is Blind:} Files within \texttt{core\_interfaces/} are strictly forbidden from importing from \texttt{infrastructure/}, \texttt{app\_*/}, or any external third-party I/O libraries (like \texttt{psycopg} or \texttt{requests}). They may only import from the Python standard library (e.g., \texttt{typing}, \texttt{dataclasses}, \texttt{abc}).
    \item \textbf{Infrastructure Conforms:} Files within \texttt{infrastructure/} \textit{must} import from \texttt{core\_interfaces/} in order to inherit and implement the required Protocols.
    \item \textbf{Applications Orchestrate:} The \texttt{app\_*} directories act as the execution boundaries. They import the pure logic from \texttt{core\_interfaces/}, instantiate the concrete adapters from \texttt{infrastructure/}, and inject the adapters into the pure logic to execute a specific use case.
\end{enumerate}

The following subsections will decompose each of these physical boundaries, detailing the exact files, classes, and responsibilities housed within them.


\textbf{The Inner Ring: \texttt{core\_interfaces/}}

The \texttt{core\_interfaces} directory is the absolute center of our system. It is computationally blind to the outside world; it contains zero external dependencies and imports only from the Python standard library. Here, we define the pure data structures (Entities) and the abstract contracts (Ports) that govern the physics of our pipeline.

We construct this inner ring sequentially, building its own internal Directed Acyclic Graph. We must define the core entity before we can define the interfaces that manipulate it.


\subsubsection*{The Domain Entity: \texttt{DocumentMetadata}}

At the very heart of our system, before we consider databases or web servers, we must define exactly what it is we are studying. For this project, our primary subject is the academic paper.

In a traditional, less rigorous software project, a developer might create a "Document" class that knows how to save itself to a hard drive or format itself for a web page. However, in our mathematically pure Core, this is strictly forbidden. Any action that interacts with the physical world—like writing to a disk or connecting via HTTP—is called a \textbf{side effect}. To maintain mathematical purity, our core components must be \textbf{pure functions} and pure data structures, meaning they produce the exact same output for a given input, every single time, without altering anything in the external environment.

Therefore, our core entity, \texttt{DocumentMetadata}, is merely a rigid container for data (a \texttt{dataclass} in Python). It holds the document's ID, title, abstract, and its high-dimensional vector representation. We freeze this structure in memory (\texttt{frozen=True}), making it immutable. It is a pure mathematical state, unaware of where it came from or where it is going.

\subsubsection*{The Port: \texttt{DocumentRepository}}

Now we encounter a fundamental problem: if our core domain is perfectly isolated and cannot interact with the physical world, how do we ever save a \texttt{DocumentMetadata} object so we don't lose our calculations when the computer turns off?

This is where we introduce the concept of a \textbf{Port} (or a Protocol).

Imagine you are designing a complex electrical circuit. You don't hardwire the circuit directly into the wall; instead, you build a standardized plug (a Port). The circuit doesn't need to know if the electricity comes from a coal plant or a solar panel; it only knows the exact shape of the plug and the voltage it expects.

In our Python code, the \texttt{DocumentRepository} is this plug. It is an abstract contract defined entirely within the \texttt{core\_interfaces} ring. It states: \textit{"I need a way to save a document, and a way to retrieve a document by its ID."} It defines the exact inputs and outputs of these functions, but it contains absolutely zero code detailing \textit{how} to actually do it. It is merely a theoretical blueprint.

\subsubsection*{The Outer Ring: Infrastructure and Adapters}

To turn our theory into physical reality, we step out into the \textbf{Infrastructure} layer (\texttt{infrastructure/database/}). This is the "outer ring," the chaotic factory floor where our program finally interacts with the operating system, network cables, and hard drives.

Here, we build the \textbf{Adapter}. If the Port is the electrical socket, the Adapter is the physical plug that connects to the specific power source.

Our adapter is a concrete Python class that implements the \texttt{Document\allowbreak{}Repository} protocol. Its job is to take our pure \texttt{DocumentMetadata} entity and translate it into a format that a database can understand.

\subsubsection*{The Need for Translation: SQLAlchemy and ORMs}

Databases do not understand Python objects. Relational databases (like PostgreSQL, which we will use) store information in rigid, two-dimensional tables made of rows and columns. They communicate using a specialized language called SQL (Structured Query Language).

To bridge the gap between our Python memory structures and the database's tabular storage, our Adapter utilizes a tool called \textbf{SQLAlchemy}. SQLAlchemy is an \textbf{Object-Relational Mapper (ORM)}. You can think of it as a highly sophisticated bilingual translator.

When the Core Domain says "save this document," the request flows outward to the Infrastructure adapter. The adapter uses SQLAlchemy to disassemble our \texttt{DocumentMetadata} object, translate its properties into a SQL \texttt{INSERT} command, open a network connection to the database, and physically write the bytes to the hard drive.

Through this Hexagonal Architecture, the mathematical purity of the Inner Ring is perfectly preserved. The Core Domain orchestrates the logic, entirely insulated from the complexities of SQL, network latency, or hard drive mechanics. We will explore the deep mechanics of Python memory management and Relational Algebra powering these outer infrastructure layers in Chapter 3 and Chapter 4.

\subsubsection*{The Machine Learning Pipeline: Hardware-Accelerated Semantics}

Having established how the Hexagonal Architecture preserves domain purity and manages persistent storage, we must address the computational core of the Semantic Document Engine: mapping raw text onto a measurable topological space. This requires orchestrating two distinct hardware accelerators: GPUs for tensor generation and CPU SIMD instructions for vector retrieval.

\paragraph{1. Asynchronous Ingestion (The Celery Workers)}
The document lifecycle begins at the application's entry point, \texttt{main.py}. Fetching thousands of documents from the arXiv API is fundamentally an I/O-bound operation; the CPU spends the majority of its clock cycles idle, waiting for TCP packets. Executing this sequentially within a single Python thread would be catastrophically slow due to the Global Interpreter Lock (GIL).

Instead, \texttt{main.py} acts as an orchestrator, delegating the ingestion workload to a distributed queue managed by \textbf{Celery}. Celery spawns independent, isolated Linux processes. By utilizing multiple processes, we completely bypass the GIL, allowing the OS scheduler to execute concurrent network syscalls and parse incoming XML payloads in parallel across all available CPU cores.

\paragraph{2. Tensor Allocation (PyTorch and VRAM)}
Once a Celery worker constructs a pure, text-based \texttt{DocumentMetadata} entity, it must compute its semantic representation. The ML Adapter invokes a Transformer model (e.g., SciBERT) via Hugging Face and PyTorch.

Physically, this operation serializes the tokenized text and transfers it across the motherboard's PCIe bus from the host's RAM into the GPU's localized VRAM. The GPU's streaming multiprocessors perform highly parallel matrix multiplications to compute the attention weights, projecting the text into a dense vector embedding $\vec{v} \in \mathbb{R}^{d}$ (where $d=768$). The resulting tensor is copied back to the host process, attached to the Domain Entity, and passed to the \texttt{DocumentRepository} for atomic insertion.

\paragraph{3. Push-Down Compute via \texttt{pgvector}}
The architectural elegance of the system emerges during a semantic search query. To find papers conceptually proximate to a user's query, we must compute the inner product (or cosine distance) between the query vector and millions of stored document vectors.

A naive architecture would query the database, pull millions of $768$-dimensional vectors across the network into Python's local heap, and compute the distances using NumPy. This design is fatally flawed: it saturates network bandwidth and triggers massive Reference Counting and Garbage Collection overhead in the CPython interpreter.

Instead, we enforce \textbf{Compute Push-Down}. By utilizing the \texttt{pgvector} C-extension in PostgreSQL, we push the mathematical operation down to the data. SQLAlchemy passes the single query vector into the SQL Abstract Syntax Tree. The PostgreSQL process executes the dot products locally using CPU SIMD (Single Instruction, Multiple Data) hardware registers. SIMD allows the processor to multiply and accumulate multiple vector dimensions in a single clock cycle, operating entirely within the CPU's hyper-fast L1/L2 cache.

PostgreSQL computes the distances, sorts the manifold, and returns only the top-$k$ nearest neighbors across the network socket. By moving the compute to the data, we reduce the network payload from gigabytes to mere kilobytes. The rigorous memory mechanics and signal processing theory underlying this pipeline will be formalized in Chapter 4 and Chapter 6.

